% Native transcription of Section 3-3, source PDF pp. 118--128
% (printed pp. 106--116). The source ends mid-sentence on PDF p. 128;
% that incomplete sentence is retained here.

% The printed Remark in this section is unnumbered; keep it out of the
% theorem counter used for the source's numbered results.
\newtheorem*{pctremark}{Remark}

% PCT-SOURCE: pdf=118 print=106 kind=prose
As was already mentioned in the introduction, it is a central problem of the
theory at the moment to find examples which are non-trivial in this sense.

% PCT-SOURCE: pdf=118 print=106 kind=section id=3-3
\subsection{Properties of the Vacuum Expectation Values}
\label{sec:pct-3-3}

% PCT-SOURCE: pdf=118 print=106 kind=prose
The object of this section is to discuss the quantities

% PCT-SOURCE: pdf=118 print=106 kind=display id=3-19
\begin{equation}
 \bra{\Omega}\varphi_1(x_1)\varphi_2(x_2)\cdots
 \varphi_n(x_n)\ket{\Omega}.
 \tag{3-19}\label{eq:pct-3-19}
\end{equation}

% PCT-SOURCE: pdf=118 print=106 kind=prose
where $\varphi_j$, $j=1,\ldots,n$, is a component of an irreducible tensor.
We shall refer to such expressions as \emph{vacuum expectation values}. Two
of the $\varphi_j$ might be components of different fields or could be the
same component of the same field or could be hermitian conjugate to each
other. In other words, we consider all such quantities, for all combinations
of the arbitrary labels, and all permutations of the indices.

% PCT-SOURCE: pdf=118 print=106 kind=prose
To give a meaning to such a symbol, we note that axiom I implies that

% PCT-SOURCE: pdf=118 print=106 kind=display id=3-20
\begin{equation}
 \bra{\Omega}\varphi_1(f_1)\varphi_2(f_2)\cdots
 \varphi_n(f_n)\ket{\Omega}
 \tag{3-20}\label{eq:pct-3-20}
\end{equation}

% PCT-SOURCE: pdf=118 print=106 kind=prose
exists and is a separately continuous multilinear functional of the
arguments $f_1,\ldots,f_n$ as they vary over $\Schwartz(\R^4)$. It follows
from the nuclear theorem (Theorem 2-1) that this functional can be uniquely
extended to be a tempered distribution of the $n$ four-vectors
$x_1,\ldots,x_n$. It is this distribution that (3-19) symbolizes, and which,
for the particular choice of field-components $\varphi_1,\ldots,\varphi_n$
(from among all components) under discussion at a given moment, we shall
denote by $\mathcal{W}$:

% PCT-SOURCE: pdf=118 print=106 kind=display id=3-21
\begin{equation}
 \mathcal{W}(x_1,x_2,\ldots,x_n)
 =\bra{\Omega}\varphi_1(x_1)\varphi_2(x_2)\cdots
 \varphi_n(x_n)\ket{\Omega}.
 \tag{3-21}\label{eq:pct-3-21}
\end{equation}

% PCT-SOURCE: pdf=118 print=106 kind=prose
With the same choice of labels we shall denote by $\mathcal{W}_\pi$ the
“permuted” vacuum expectation value

% PCT-SOURCE: pdf=118 print=106 kind=display id=3-22
\begin{equation}
 \mathcal{W}_\pi(x_1,x_2,\ldots,x_n)
 =\bra{\Omega}\varphi_{i_1}(x_{i_1})\varphi_{i_2}(x_{i_2})\cdots
 \varphi_{i_n}(x_{i_n})\ket{\Omega},
 \tag{3-22}\label{eq:pct-3-22}
\end{equation}

% PCT-SOURCE: pdf=118 print=106 kind=prose
where $\pi$ is the permutation
$(1,2,\ldots,n)\longmapsto(i_1,i_2,\ldots,i_n)$. Clearly if all the
components $\varphi_j$ are different there are $n!$ different $\mathcal{W}_\pi$
for each choice of components, and fewer if some coincide.

% PCT-SOURCE: pdf=118 print=106 kind=prose
The argument used above to give a meaning to (3-19) as a distribution can
also be used to make sense of expressions of the form

% PCT-SOURCE: pdf=118 print=106 kind=display id=3-23
\begin{equation}
 \ket{\Psi}=\int \dd^4x_1\cdots\dd^4x_k\,
 f(x_1,\ldots,x_k)\varphi_1(x_1)\cdots\varphi_k(x_k)\ket{\Omega}.
 \tag{3-23}\label{eq:pct-3-23}
\end{equation}

% PCT-SOURCE: pdf=118 print=106 kind=prose
We can find a sequence of functions $\{f_J\}$, with
% PCT-SOURCE: pdf=118 print=106 kind=display
\begin{equation*}
 f_J(x_1,\ldots,x_k)=\sum_{j=1}^{J}f_1^j(x_1)\cdots f_k^j(x_k).
\end{equation*}

% PCT-SOURCE: pdf=119 print=107 kind=prose
and $f_j^i\in\Schwartz(\R^4)$ such that $f_J\to f$ in
$\Schwartz(\R^{4k})$ as $J\to\infty$. Then the sequence of states

% PCT-SOURCE: pdf=119 print=107 kind=display id=3-24
\begin{equation}
 \ket{\Psi_J}=\sum_{j=1}^{J}
 \varphi_1(f_1^j)\varphi_2(f_2^j)\cdots\varphi_k(f_k^j)\ket{\Omega}
 \tag{3-24}\label{eq:pct-3-24}
\end{equation}

% PCT-SOURCE: pdf=119 print=107 kind=prose
converges in norm, and we take its limit as the definition of $\ket{\Psi}$
in (3-23). To prove the convergence of (3-24), consider

% PCT-SOURCE: pdf=119 print=107 kind=display
\begin{equation*}
 \begin{gathered}
  \left\|\ket{\Psi_m}-\ket{\Psi_n}\right\|^2\\[-2pt]
  =\left\langle
  \begin{gathered}
   \sum_{j=m}^{n}\\[-2pt]
   {}\varphi_1(f_1^j)\cdots\\[-2pt]
   {}\varphi_k(f_k^j)\ket{\Omega}
  \end{gathered}
  \middle|
  \begin{gathered}
   \sum_{j=m}^{n}\\[-2pt]
   {}\varphi_1(f_1^j)\cdots\\[-2pt]
   {}\varphi_k(f_k^j)\ket{\Omega}
  \end{gathered}
  \right\rangle\\
  \begin{aligned}
   ={ }&\int \mathcal{W}(y_k,\ldots,y_1,x_1,\ldots,x_k)\,
   \overline{(f_n-f_m)(y_1,\ldots,y_k)}\\[-2pt]
   &\qquad\times(f_n-f_m)(x_1,\ldots,x_k)\\[-2pt]
   &\qquad\quad\times\dd^4x_1\cdots\dd^4x_k\\[-2pt]
   &\qquad\quad\times\dd^4y_1\cdots\dd^4y_k,
  \end{aligned}
 \end{gathered}
\end{equation*}

% PCT-SOURCE: pdf=119 print=107 kind=prose
where $\mathcal{W}$ corresponds to the fields
$\varphi_k^\dagger\cdots\varphi_1^\dagger\varphi_1\cdots\varphi_k$ in that order.
Now the function $\overline{(f_n-f_m)}(f_n-f_m)$ tends to zero in
$\Schwartz$, and since $\mathcal{W}$ is a tempered distribution,
$\left\|\ket{\Psi_m}-\ket{\Psi_n}\right\|\to0$. Since $\mathcal{H}$ is a
complete space,\footnote{See Section~2-6.} $\{\ket{\Psi_m}\}$ has a limit
state $\ket{\Psi}$.

% PCT-SOURCE: pdf=119 print=107 kind=prose
The set of all vectors symbolized by (3-23) will be denoted $D_1$. Clearly,
any $\varphi_j(f)$ can be defined by continuity on $D_1$,

% PCT-SOURCE: pdf=119 print=107 kind=display
\begin{equation*}
 \begin{aligned}
 \varphi_j(f)\ket{\Psi}
 &=\int \dd^4x_1\cdots\dd^4x_k\,\dd^4x\,
 f(x)f(x_1,\ldots,x_k)\varphi_j(x)\\[-2pt]
 &\qquad\times\varphi_1(x_1)\cdots\varphi_k(x_k)\ket{\Omega}.
 \end{aligned}
\end{equation*}

% PCT-SOURCE: pdf=119 print=107 kind=prose
It is easy to verify that the knowledge of the $\varphi_j$ on the domain
$D_0$ and the validity of I, II, and III there imply I, II, and III on $D_1$.

% PCT-SOURCE: pdf=119 print=107 kind=prose
We express our assertions about the vacuum expectation values in several
theorems. The first is

% PCT-SOURCE: pdf=119 print=107 kind=theorem id=3-1
\begin{theorem}
\textbf{(a) (\emph{Relativistic Transformation Law}).}\par
Suppose $\varphi_\alpha,\psi_\beta,\ldots,\chi_\gamma$ are $n$ fields
which transform according to irreducible representations of
$\mathrm{SL}(2,\C)$, namely

% PCT-SOURCE: pdf=119 print=107 kind=display id=3-25
\begin{equation}
 U(a,A)\varphi_\alpha(x)U(a,A)^{-1}
 =\sum_{\alpha'}S^{(\varphi)}_{\alpha\alpha'}(A^{-1})
 \varphi_{\alpha'}(Ax+a).
 \tag{3-25}\label{eq:pct-3-25}
\end{equation}

% PCT-SOURCE: pdf=119 print=107 kind=prose
[$\psi,\ldots,\chi$ transforming with
$S^{(\psi)},\ldots,S^{(\chi)}$, respectively.] Then the vacuum
expectation values for $n=1,2,3,\ldots$ are tempered distributions which
transform according to

% PCT-SOURCE: pdf=119 print=107 kind=display id=3-26
\begin{align}
 &\sum_{\alpha',\beta',\ldots,\gamma'}
 S^{(\varphi)}_{\alpha\alpha'}(A)S^{(\psi)}_{\beta\beta'}(A)\cdots
 S^{(\chi)}_{\gamma\gamma'}(A)
 \bra{\Omega}\varphi_{\alpha'}(x_1)\psi_{\beta'}(x_2)\cdots
 \chi_{\gamma'}(x_n)\ket{\Omega}
 \notag\\
 &\qquad=\bra{\Omega}\varphi_\alpha(Ax_1+a)\psi_\beta(Ax_2+a)\cdots
 \chi_\gamma(Ax_n+a)\ket{\Omega}.
 \tag{3-26}\label{eq:pct-3-26}
\end{align}
\end{theorem}

% PCT-SOURCE: pdf=120 print=108 kind=proof id=3-1
\begin{proof}
As remarked earlier, the vacuum expectation values are tempered
distributions, by the nuclear theorem. Using (3-25) with $A$ replaced by
$A^{-1}$, and the property $U(a,A)\ket{\Omega}=\ket{\Omega}$ [Eq.~(3-1)],
(3-26) follows immediately.
\end{proof}

% PCT-SOURCE: pdf=120 print=108 kind=theorem id=3-2
\begin{theorem}
Suppose $\varphi_1,\varphi_2,\ldots,\varphi_n$ are any components of any
fields, and $\mathcal{W}$ is defined as in (3-21). Then

% PCT-SOURCE: pdf=120 print=108 kind=prose
\textbf{(b) (\emph{Spectral Conditions}).}\par
There are tempered distributions
$W(\xi_1,\ldots,\xi_{n-1})$ depending on the relative coordinates

% PCT-SOURCE: pdf=120 print=108 kind=display id=3-27
\begin{equation}
 \xi_j=x_j-x_{j+1},\qquad j=1,2,\ldots,n-1
 \tag{3-27}\label{eq:pct-3-27}
\end{equation}

% PCT-SOURCE: pdf=120 print=108 kind=prose
which satisfy

% PCT-SOURCE: pdf=120 print=108 kind=display id=3-28
\begin{equation}
 \mathcal{W}(x_1,\ldots,x_n)=W(\xi_1,\xi_2,\ldots,\xi_{n-1}).
 \tag{3-28}\label{eq:pct-3-28}
\end{equation}

% PCT-SOURCE: pdf=120 print=108 kind=prose
The Fourier transforms of $\mathcal{W}$ and $W$ are tempered distributions
defined by

% PCT-SOURCE: pdf=120 print=108 kind=display id=3-29
\begin{equation}
 \widetilde{\mathcal{W}}(p_1,p_2,\ldots,p_n)
 =\int \exp\!\left(-\ii\sum_{j=1}^{n}p_j\!\cdot x_j\right)
 \mathcal{W}(x_1,\ldots,x_n)\,\dd^4x_1\cdots\dd^4x_n
 \tag{3-29}\label{eq:pct-3-29}
\end{equation}

% PCT-SOURCE: pdf=120 print=108 kind=display id=3-30
\begin{equation}
 \widetilde{W}(q_1,\ldots,q_{n-1})
 =\int \exp\!\left(-\ii\sum_{j=1}^{n-1}q_j\!\cdot\xi_j\right)
 W(\xi_1,\xi_2,\ldots,\xi_{n-1})\,\dd^4\xi_1\cdots\dd^4\xi_{n-1}
 \tag{3-30}\label{eq:pct-3-30}
\end{equation}

% PCT-SOURCE: pdf=120 print=108 kind=prose
and are related by

% PCT-SOURCE: pdf=120 print=108 kind=display id=3-31
\begin{equation}
 \widetilde{\mathcal{W}}(p_1,\ldots,p_n)
 =(2\pi)^4\delta^{(4)}\!\left(\sum_{j=1}^{n}p_j\right)
 \widetilde{W}(p_1,p_1+p_2,\ldots,p_1+p_2+\cdots+p_{n-1}).
 \tag{3-31}\label{eq:pct-3-31}
\end{equation}

% PCT-SOURCE: pdf=120 print=108 kind=prose
Further we have

% PCT-SOURCE: pdf=120 print=108 kind=display id=3-32
\begin{equation}
 \widetilde{W}(q_1,\ldots,q_{n-1})=0
 \tag{3-32}\label{eq:pct-3-32}
\end{equation}

% PCT-SOURCE: pdf=120 print=108 kind=prose
if any $q$ is not in the energy momentum spectrum of the states.

% PCT-SOURCE: pdf=120 print=108 kind=prose
\textbf{(c) (\emph{Hermiticity Conditions})}

% PCT-SOURCE: pdf=120 print=108 kind=display id=3-33
\begin{equation}
 \bra{\Omega}\varphi_1(x_1)\varphi_2(x_2)\cdots\varphi_n(x_n)\ket{\Omega}
 =\overline{\bra{\Omega}\varphi_n^\dagger(x_n)\cdots
 \varphi_2^\dagger(x_2)\varphi_1^\dagger(x_1)\ket{\Omega}}.
 \tag{3-33}\label{eq:pct-3-33}
\end{equation}

% PCT-SOURCE: pdf=121 print=109 kind=prose
\textbf{(d) (\emph{Local Commutativity Conditions})}

% PCT-SOURCE: pdf=121 print=109 kind=display id=3-34
\begin{equation}
 \begin{aligned}
 \mathcal{W}(x_1,x_2,\ldots,x_{j+1},x_j,\ldots,x_n)
 &=(-1)^m\mathcal{W}(x_1,x_2,\ldots,x_j,x_{j+1},\ldots,x_n),\\
 &\qquad j=1,\ldots,n-1
 \end{aligned}
 \tag{3-34}\label{eq:pct-3-34}
\end{equation}

% PCT-SOURCE: pdf=121 print=109 kind=prose
if the difference $x_j-x_{j+1}$ is space-like, where $m$ is zero if
$\varphi_j$ and $\varphi_{j+1}$ commute, and one if they anti-commute.
\end{theorem}

% PCT-SOURCE: pdf=121 print=109 kind=proof id=3-2
\begin{proof}
The argument given in Section 2-1 shows that translation invariance implies
that $\mathcal{W}$ depends only on the differences
$\xi_1,\ldots,\xi_{n-1}$, or more precisely, that there exists a tempered
distribution $W$ satisfying (3-28). Clearly

% PCT-SOURCE: pdf=121 print=109 kind=display
\begin{align*}
 &\int \dd^4x_1\cdots\dd^4x_n\,
 \exp\!\left(-\ii\sum_{j=1}^{n}p_j\!\cdot x_j\right)
 \mathcal{W}(x_1,x_2,\ldots,x_n) \\
 &\quad=\int \dd^4x_1\cdots\dd^4x_n\,
 \exp\!\bigl\{-\ii[p_1\!\cdot(x_1-x_2)
 +(p_1+p_2)\!\cdot(x_2-x_3)+\cdots \\
 &\hspace{42mm}+(p_1+\cdots+p_{n-1})\!\cdot(x_{n-1}-x_n)]\bigr\} \\
 &\qquad\quad\times\exp\!\left[-\ii\left(\sum_{j=1}^{n}p_j\right)\!\cdot x_n\right]
 W(x_1-x_2,\ldots,x_{n-1}-x_n) \\
 &\quad=(2\pi)^4\delta^{(4)}\!\left(\sum_{j=1}^{n}p_j\right)
 \widetilde{W}(p_1,p_1+p_2,\ldots,p_1+p_2+\cdots+p_{n-1}).
\end{align*}

% PCT-SOURCE: pdf=121 print=109 kind=prose
As explained in Section 2-6, for any two states $\ket{\Phi}$, $\ket{\Psi}$,
we have

% PCT-SOURCE: pdf=121 print=109 kind=display
\begin{equation*}
 \int \exp(+\ii p\!\cdot a)\,\dd^4a\,
 \bra{\Phi}U(a,1)\ket{\Psi}=0,
\end{equation*}

% PCT-SOURCE: pdf=121 print=109 kind=prose
unless $p$ lies in the energy-momentum spectrum of the states. Therefore

% PCT-SOURCE: pdf=121 print=109 kind=display
\begin{equation*}
 \int \exp(-\ii p\!\cdot a)\,\dd^4a\,
 \bra{\Omega}\varphi_1(f_1)\cdots\varphi_j(f_j)U(-a,1)
 \varphi_{j+1}(f_{j+1})\cdots\varphi_n(f_n)\ket{\Omega}=0,
\end{equation*}

% PCT-SOURCE: pdf=121 print=109 kind=prose
which implies

% PCT-SOURCE: pdf=121 print=109 kind=display
\begin{equation*}
 \begin{aligned}
  \int \exp(-\ii p\!\cdot a)\,
  W(\xi_1,\ldots,\xi_{j-1},\xi_j+a,\xi_{j+1},\ldots,\xi_{n-1})\,\dd^4a&=0,\\[-2pt]
  &\qquad j=1,2,\ldots,n
 \end{aligned}
\end{equation*}

% PCT-SOURCE: pdf=121 print=109 kind=prose
for $p$ not in the physical spectrum, that is,
$\widetilde{W}(q_1,\ldots,q_{n-1})=0$ unless each $q_j$ lies in the physical
spectrum.

% PCT-SOURCE: pdf=121 print=109 kind=prose
\textbf{(c)} This relation follows immediately from

% PCT-SOURCE: pdf=121 print=109 kind=display
\begin{equation*}
 \bra{\Omega}\varphi_1(f_1)\cdots\varphi_n(f_n)\ket{\Omega}
 =\overline{\bra{\Omega}\varphi_n^\dagger(f_n)\cdots
 \varphi_1^\dagger(f_1)\ket{\Omega}}
\end{equation*}

% PCT-SOURCE: pdf=121 print=109 kind=prose
and the fact that we can continue this relation from test functions of the
form $f_1(x_1)\allowbreak f_2(x_2)\allowbreak\cdots\allowbreak f_n(x_n)$ to
all of $\Schwartz(\R^{4n})$.

% PCT-SOURCE: pdf=122 print=110 kind=prose
\textbf{(d)} This relation follows immediately from axiom III, again with the
use of the extension from product test functions to all test functions of
$\Schwartz(\R^{4n})$.
\end{proof}

% PCT-SOURCE: pdf=122 print=110 kind=prose
The next property follows from the positive definiteness of the scalar
product in Hilbert space.

% PCT-SOURCE: pdf=122 print=110 kind=theorem id=3-3
\begin{theorem}
\textbf{(e) (\emph{Positive Definiteness Conditions}).} For any sequence
$\{f_j\}$ of test functions, $f_j\in\Schwartz(\R^{4j})$, with $f_j=0$
except for a finite number of $j$, the vacuum expectation values satisfy the
inequalities

% PCT-SOURCE: pdf=122 print=110 kind=display id=3-35
\begin{equation}
 \begin{aligned}
 \sum_{j,k=0}^{\infty}\int\cdots\int
 &\overline{f_j(x_1,\ldots,x_j)}\,
 \mathcal{W}_{jk}(x_j,x_{j-1},\ldots,x_1,y_1,\ldots,y_k)
 f_k(y_1,\ldots,y_k)\\
 &\qquad\times\dd^4x_1\cdots\dd^4x_j\,\dd^4y_1\cdots\dd^4y_k\geq0.
 \end{aligned}
 \tag{3-35}\label{eq:pct-3-35}
\end{equation}

% PCT-SOURCE: pdf=122 print=110 kind=prose
In (3-35) $\mathcal{W}_{jk}$ represents

% PCT-SOURCE: pdf=122 print=110 kind=display
\begin{equation*}
 \bra{\Omega}\varphi_{jj}^{\dagger}(x_j)\cdots
 \varphi_{j1}^{\dagger}(x_1)\varphi_{k1}(y_1)\cdots
 \varphi_{kk}(y_k)\ket{\Omega},
\end{equation*}

% PCT-SOURCE: pdf=122 print=110 kind=prose
where the field components labeled by $jk$ can be any selection from among
all components in the theory. Moreover, if, for any sequence
$f_0,f_1,\ldots$ we have the equality sign in (3-35), then\footnote{It is
shown on p.~121 that this follows from (3-35) and the other axioms.} it must
give 0 also when the sequence, $\{f_j\}$, is replaced by a sequence
$\{g_j\}$ where

% PCT-SOURCE: pdf=122 print=110 kind=display id=3-36
\begin{equation}
 \begin{aligned}
  g_0&=0,\qquad g_1=g(x_1)f_0,\\
  g_2&=g(x_1)f_1(x_2),\qquad
  g_3=g(x_1)f_2(x_2,x_3)\cdots
 \end{aligned}
\tag{3-36}\label{eq:pct-3-36}
\end{equation}

% PCT-SOURCE: pdf=122 print=110 kind=prose
for any test function $g$.
\end{theorem}

% PCT-SOURCE: pdf=122 print=110 kind=proof id=3-3
\begin{proof}
The inequalities (3-35) express the fact that the norm of the state

% PCT-SOURCE: pdf=122 print=110 kind=display
\begin{align*}
 \ket{\Psi}={}&f_0\ket{\Omega}+\varphi_{11}(f_1)\ket{\Omega}
 +\int\varphi_{21}(x_1)\varphi_{22}(x_2)f_2(x_1,x_2)
 \dd^4x_1\dd^4x_2\ket{\Omega}+\cdots\\
 &+\int\varphi_{j1}(x_1)\varphi_{j2}(x_2)\cdots\varphi_{jj}(x_j)
 f_j(x_1,\ldots,x_j)\,\dd^4x_1\cdots\dd^4x_j\ket{\Omega}+\cdots
\end{align*}

% PCT-SOURCE: pdf=122 print=110 kind=prose
is non-negative. If the norm is 0, $\ket{\Psi}=0$, and so
$\varphi_\alpha(g)\ket{\Psi}=0$ for any component $\varphi_\alpha$ and any
test function $g$. Hence the expression (3-35) must give 0 also for the
sequence (3-36).
\end{proof}

% PCT-SOURCE: pdf=122 print=110 kind=prose
As we shall see in the next section, the conditions (a) through (e) are
sufficient to guarantee that a set of tempered distributions are the vacuum
expectation values of some field theory satisfying O, I, II, and III, with
the exception of the requirement that the vacuum be unique. The next theorem
gives an additional condition which ensures this property, as we shall see in
Section 3-4.

% PCT-SOURCE: pdf=123 print=111 kind=theorem id=3-4
\begin{theorem}
\textbf{(\emph{Cluster Decomposition Property}).} If $a$ is a space-like
vector, then

% PCT-SOURCE: pdf=123 print=111 kind=display id=3-37
\begin{equation}
 \begin{aligned}
  &\mathcal{W}(x_1,\ldots,x_j,x_{j+1}+\lambda a,x_{j+2}+\lambda a,\ldots,\\[-2pt]
  &\qquad x_n+\lambda a)\\[-2pt]
  &\longrightarrow
  \mathcal{W}(x_1,\ldots,x_j)\mathcal{W}(x_{j+1},\ldots,x_n)
 \end{aligned}
\tag{3-37}\label{eq:pct-3-37}
\end{equation}

% PCT-SOURCE: pdf=123 print=111 kind=prose
as $\lambda\to\infty$, in the sense of convergence in $\Schwartz'$.
\end{theorem}

% PCT-SOURCE: pdf=123 print=111 kind=remark
\begin{pctremark}
We shall prove the theorem only for a theory in which there is no state other
than the vacuum having a mass less than some positive threshold $M$ (which may
be as small as we like). Such theories will be said to have a \emph{mass
gap}. If there are mass zero particles in the theory there is no mass gap,
but (3-37) still holds; the proof is more technical in this case. Equation
(3-37) can also be proved without assuming axiom III. See Refs.~5, 7, 9, 10,
and 11 for further information.
\end{pctremark}

% PCT-SOURCE: pdf=123 print=111 kind=proof id=3-4
\begin{proof}
The proof is based on a very simple argument due to D.~Ruelle. It exploits
the fact that for large $\lambda$ the operators occurring in (3-37) can be
rewritten in the order $x_{j+1}+\lambda a,\ldots,x_n+\lambda a,x_1,\ldots,x_j$
with at most a change in sign in the function. This reversal of order has
the effect of reversing the sign of the momentum conjugate to the difference
variable $x_j-x_{j+1}$. The change of sign makes it possible to use the
spectral condition in this momentum vector to get the theorem. The details
are as follows.

% PCT-SOURCE: pdf=123 print=111 kind=prose
For convenience, choose $a^2=1$. This can be done without loss of
generality by changing the scale of $\lambda$. The Fourier transforms of the
distributions

% PCT-SOURCE: pdf=123 print=111 kind=display
\begin{align*}
 F_1={}&\mathcal{W}(x_1,\ldots,x_j,x_{j+1}+\lambda a,\ldots,x_n+\lambda a)\\[-2pt]
 &-\mathcal{W}(x_1,\ldots,x_j)\mathcal{W}(x_{j+1},\ldots,x_n),\\
 F_2={}&\mathcal{W}(x_{j+1}+\lambda a,\ldots,x_n+\lambda a,x_1,\ldots,x_j)\\[-2pt]
 &-\mathcal{W}(x_1,\ldots,x_j)\mathcal{W}(x_{j+1},\ldots,x_n)
\end{align*}

% PCT-SOURCE: pdf=123 print=111 kind=prose
are zero unless $P=p_1+\cdots+p_j$ lies in the forward and backward cones,
respectively, by property (b). If there is a unique vacuum state,
$\ket{\Omega}$, then the product
$\mathcal{W}(x_1,\ldots,x_j)\allowbreak
\mathcal{W}(x_{j+1},\ldots,x_n)$ cancels the contribution from this state at
the origin $P^\mu=0$, and we get
$\widetilde{F}_1\pm\widetilde{F}_2=0$ unless $P^2\leq -M^2$. For fixed
$x_1,\ldots,x_n$, $F_1$ and $F_2$ coincide up to a sign for large enough
$\lambda$, by locality. From Section 2-1, $F_1\pm F_2$ is a finite sum of
derivatives of a polynomially bounded continuous function, $G$ say. Thus,
for any test function $h\in\Schwartz$,

% PCT-SOURCE: pdf=123 print=111 kind=display
\begin{equation*}
 \begin{aligned}
 &\int (F_1\pm F_2)(x_1,\ldots,x_n)h(x_1,\ldots,x_n)\\[-2pt]
 &\qquad\times\dd^4x_1\cdots\dd^4x_n\\
 &={ }\int D^mG(\lambda,x_1,\ldots,x_n)h(x_1,\ldots,x_n)\\[-2pt]
 &\qquad\times\dd^4x_1\cdots\dd^4x_n.
 \end{aligned}
\end{equation*}

% PCT-SOURCE: pdf=124 print=112 kind=prose
Now $G$ is polynomially bounded in all variables, i.e.,

% PCT-SOURCE: pdf=124 print=112 kind=display
\begin{equation*}
 |G|\leq G_0\lambda^N R^Q
\end{equation*}

% PCT-SOURCE: pdf=124 print=112 kind=prose
for sufficiently large $\lambda$ and $R$, where by $R$ we mean the
Euclidean norm in $\R^{4n}$ over which the integral runs, i.e.,

% PCT-SOURCE: pdf=124 print=112 kind=display
\begin{equation*}
 R^2=\sum_{i=1}^{n}\bigl[(x_i^0)^2+\mathbf{x}_i^2\bigr].
\end{equation*}

% PCT-SOURCE: pdf=124 print=112 kind=prose
Using local commutativity we can and will choose the $\pm$ sign so that
$D^mG=0$ for large $\lambda$. We shall need a more quantitative form of
this statement: $D^mG=0$ for $R<R_0$, where $R_0$ is a positive multiple of
$\lambda$. To determine $R_0$ note that

% PCT-SOURCE: pdf=124 print=112 kind=display
\begin{align*}
 (x_i-x_k-\lambda a)^2
 &=-(x_i^0-x_k^0)^2+(\mathbf{x}_i-\mathbf{x}_k)^2+\lambda^2\\[-2pt]
 &\quad+2\lambda\bigl(a^0(x_i^0-x_k^0)-\mathbf{a}\!\cdot(\mathbf{x}_i-\mathbf{x}_k)\bigr)\\
 &>-|x_i^0|^2-|x_k^0|^2-2|x_i^0|\,|x_k^0|+\lambda^2\\
 &\qquad-2\lambda\bigl(|a^0|(|x_i^0|+|x_k^0|)
 +|\mathbf{a}|(|\mathbf{x}_i|+|\mathbf{x}_k|)\bigr),
\end{align*}

% PCT-SOURCE: pdf=124 print=112 kind=prose
so that when

% PCT-SOURCE: pdf=124 print=112 kind=display
\begin{equation*}
 R^2=\sum_{i=1}^{n}\bigl[(x_i^0)^2+\mathbf{x}_i^2\bigr]<R_0^2
\end{equation*}

% PCT-SOURCE: pdf=124 print=112 kind=prose
we have

% PCT-SOURCE: pdf=124 print=112 kind=display
\begin{equation*}
 (x_i-x_k-\lambda a)^2>\lambda^2-4R_0^2
 -4\lambda R_0(|a^0|+|\mathbf{a}|).
\end{equation*}

% PCT-SOURCE: pdf=124 print=112 kind=prose
The right-hand side is $>0$ for all positive $\lambda$ if $R_0$ is chosen as
a sufficiently small multiple of $\lambda$, say
$R_0=\frac18[|a^0|+|\mathbf{a}|]^{-1}\lambda$. From this estimate follows,
for any $\epsilon>0$ and sufficiently small,

% PCT-SOURCE: pdf=124 print=112 kind=display
\begin{align*}
 \left|\begin{gathered}
  \int D^mG(\lambda,x_1,\ldots,x_n)h(x_1,\ldots,x_n)\\[-2pt]
  {}\times\dd^4x_1\cdots\dd^4x_n
 \end{gathered}\right|
 &=\left|\begin{gathered}
  \int_{R>R_0-\epsilon}D^mG(\lambda,x_1,\ldots,x_n)\\[-2pt]
  {}\times h(x_1,\ldots,x_n)\\[-2pt]
  {}\times\dd^4x_1\cdots\dd^4x_n
 \end{gathered}\right|\\
 &=\left|\begin{gathered}
  \int_{R>R_0-\epsilon}G(\lambda,x_1,\ldots,x_n)\\[-2pt]
  {}\times D^mh(x_1,\ldots,x_n)\\[-2pt]
  {}\times\dd^4x_1\cdots\dd^4x_n
 \end{gathered}\right|\\
 &\leq\int_{R>R_0-\epsilon}|G|\,|D^mh|\,R^{4n-1}\,\dd R\,\dd\omega.
\end{align*}

% PCT-SOURCE: pdf=124 print=112 kind=prose
where in the last expression we have passed to polar coordinates. Since for
any $q$, the inequality $|D^mh|<cR^{-q}$ holds for some $c$ and all
sufficiently large $R$, the right-hand side of the preceding inequality is
bounded by

% PCT-SOURCE: pdf=124 print=112 kind=display
\begin{equation*}
 \int_{R>R_0}G_0\lambda^N R^{Q-q+4n-1}\,\dd R\,\dd\omega
\end{equation*}

% PCT-SOURCE: pdf=125 print=113 kind=prose
for all sufficiently large $\lambda$. For sufficiently large $q$ this
integral converges and shows that

% PCT-SOURCE: pdf=125 print=113 kind=display
\begin{equation*}
 \left|\int(D^mG)h\right|<\frac{c_1}{\lambda^p}
\end{equation*}

% PCT-SOURCE: pdf=125 print=113 kind=prose
for all $p$, some $c_1$, and all sufficiently large $\lambda$. This is what
is meant by

% PCT-SOURCE: pdf=125 print=113 kind=display
\begin{equation*}
 \lambda^p(F_1\pm F_2)(x_1,\ldots,x_j,x_{j+1}+\lambda a,\ldots,
 x_n+\lambda a)\longrightarrow0
\end{equation*}

% PCT-SOURCE: pdf=125 print=113 kind=prose
as $\lambda\to\infty$ in the sense of convergence in $\Schwartz'$.

% PCT-SOURCE: pdf=125 print=113 kind=prose
The final step of the proof is to show that this equation holds for $F_1$ and
$F_2$ separately, and here is where the spectral condition is used. Replace
the $h$ of the preceding argument by a new $h_1$, which is defined by
$\widetilde h_1=\theta\widetilde h$ where $\theta$ is an infinitely
differentiable function of the variable $P=\sum_{k=1}^{j}p_k$, equal to 1
 for $P^2\leq -M^2$, $P^0>0$, and 0 for $P^0\leq0$. Clearly
$h_1\in\Schwartz$ so the preceding argument works equally well for it. But
$\int F_2h_1\,\dd^4x_1\cdots\dd^4x_n=0$ and
$\int F_1h_1\,\dd^4x_1\cdots\dd^4x_n=\int F_1h\,\dd^4x_1\cdots\dd^4x_n$,
so the required result (3-37) follows.
\end{proof}

% PCT-SOURCE: pdf=125 print=113 kind=prose
The physical meaning of the cluster decomposition property is that, when two
systems at points $x$ and $y$ become separated by a large space-like distance,
the interaction between them falls off to zero. The method of proof given
here shows that if there is a mass gap, the limit converges faster than any
inverse power $1/\lambda$. It can be shown to decrease exponentially, the
damping factor depending on the threshold mass $M$. If there are zero-mass
particles in the theory the limit may go as slowly as $1/\lambda^2$; this is
just the Coulomb force!

% PCT-SOURCE: pdf=125 print=113 kind=prose
If the fields have doubled components so that (1-52) makes sense, we can
formulate the requirement that $U(I_s)$ and $U(C)$ are unitary operators by
substituting these transformation laws on each component of the product of
fields entering in the state (3-23), and requiring that all scalar products be
left invariant. This leads to further properties of the $\mathcal{W}$'s. For
example, if charge conjugation defines a symmetry, then

% PCT-SOURCE: pdf=125 print=113 kind=display id=3-38
\begin{equation}
 \bra{\Omega}\varphi_{(\alpha)(\beta)}(x_1)\cdots
 \psi_{(\mu)(\nu)}(x_n)\ket{\Omega}
 =\bra{\Omega}\varphi^*_{(\dot\alpha)(\dot\beta)}(x_1)\cdots
 \psi^*_{(\dot\mu)(\dot\nu)}(x_n)\ket{\Omega}
 \tag{3-38}\label{eq:pct-3-38}
\end{equation}

% PCT-SOURCE: pdf=125 print=113 kind=prose
holds for all products in all orders. Equation (3-38) is a further
restriction on the $\mathcal{W}$'s, and should not be confused with the rule
for taking the hermitian conjugate, (3-33). If $\Theta$ is a symmetry of the
theory we get the relations

% PCT-SOURCE: pdf=125 print=113 kind=display id=3-39
\begin{equation}
 \begin{aligned}
  \bra{\Omega}\varphi_{(\alpha)(\beta)}(x_1)\cdots
  \psi_{(\mu)(\nu)}(x_n)\ket{\Omega}
  &=\ii^F(-1)^J\\[-2pt]
  &\quad\times\overline{\bra{\Omega}\varphi^*_{(\alpha)(\beta)}(-x_1)\cdots
  \psi^*_{(\mu)(\nu)}(-x_n)\ket{\Omega}}
 \end{aligned}
\tag{3-39}\label{eq:pct-3-39}
\end{equation}

% PCT-SOURCE: pdf=125 print=113 kind=prose
from (1-53), since $\Theta$ is anti-unitary, where $F$ is the total number
of half-odd integer spin fields and $J$ is the total number of undotted
indices in the collections $(\alpha),\ldots,(\mu)$.

% PCT-SOURCE: pdf=126 print=114 kind=prose
In the same way, equations similar to (3-38) and (3-39) can be found if
$I_s$ and $I_t$ define symmetries.

% PCT-SOURCE: pdf=126 print=114 kind=prose
We now introduce an important technique in the study of local fields; we
relate the vacuum expectation values to holomorphic functions.

% PCT-SOURCE: pdf=126 print=114 kind=theorem id=3-5
\begin{theorem}
The $\mathcal{W}$ and $W$ are boundary values of holomorphic functions in the
following sense. There exist holomorphic functions $\mathcal{W}$ and $W$,
such that

% PCT-SOURCE: pdf=126 print=114 kind=display
\begin{equation*}
 \mathcal{W}(z_1,\ldots,z_n)
 =W(z_1-z_2,z_2-z_3,\ldots,z_{n-1}-z_n),
\end{equation*}

% PCT-SOURCE: pdf=126 print=114 kind=prose
and $W$ is holomorphic in the variables
$\zeta_j=\xi_j-\ii\eta_j=z_j-z_{j+1}$, $j=1,\ldots,n$, the domain of
holomorphy being the tube
$\{\xi_1,\ldots,\xi_{n-1}\mid\eta_j\in V_+\}$, and

% PCT-SOURCE: pdf=126 print=114 kind=display
\begin{equation*}
 W(\xi_1,\ldots,\xi_{n-1})
 =\lim_{\eta_1,\ldots,\eta_{n-1}\to0}
 W(\xi_1-\ii\eta_1,\ldots,\xi_{n-1}-\ii\eta_{n-1}),
\end{equation*}

% PCT-SOURCE: pdf=126 print=114 kind=prose
where the limit is to be understood in the sense of convergence in
$\Schwartz'$. Furthermore, $W$ is polynomially bounded\footnote{Defined in
Section 2-3.} in $\Tube_{n-1}$. The holomorphic functions $W$ possess a
single-valued continuation into the extended tube $\ExtendedTube_{n-1}$.
\end{theorem}

% PCT-SOURCE: pdf=126 print=114 kind=proof id=3-5
\begin{proof}
Since $\widetilde{W}(q_1,\ldots,q_{n-1})=0$ unless each $q$ lies in the
forward cone, the Laplace transform

% PCT-SOURCE: pdf=126 print=114 kind=display
\begin{equation*}
 \begin{aligned}
 W(\xi_1-\ii\eta_1,\ldots,\xi_{n-1}-\ii\eta_{n-1})
 ={}&(2\pi)^{-4(n-1)}\int\cdots\int\\[-2pt]
 &\quad\dd^4q_1\cdots\dd^4q_{n-1}\\[-2pt]
 &\quad\times\exp\!\left[+\ii\sum_{j=1}^{n-1}
 (\xi_j-\ii\eta_j)\!\cdot q_j\right]\\[-2pt]
 &\qquad\widetilde{W}(q_1,\ldots,q_{n-1})
 \end{aligned}
\end{equation*}

% PCT-SOURCE: pdf=126 print=114 kind=prose
is a holomorphic function, by Theorem 2-8. The other assertions, except for
the last, follow from the same theorem.

% PCT-SOURCE: pdf=126 print=114 kind=prose
We now make use of Lorentz invariance to show that each vacuum expectation
value has a one-valued continuation into a much larger region, the extended
tube $\ExtendedTube_{n-1}$. Consider (3-26):

% PCT-SOURCE: pdf=126 print=114 kind=display
\begin{align*}
 &\sum_{\alpha',\beta',\ldots,\gamma'}
 S^{(\varphi)}_{\alpha\alpha'}(A)S^{(\psi)}_{\beta\beta'}(A)\cdots
 S^{(\chi)}_{\gamma\gamma'}(A)
 \bra{\Omega}\varphi_{\alpha'}(x_1)\psi_{\beta'}(x_2)\cdots
 \chi_{\gamma'}(x_n)\ket{\Omega} \\
 &\qquad=\bra{\Omega}\varphi_\alpha(Ax_1+a)\psi_\beta(Ax_2+a)\cdots
 \chi_\gamma(Ax_n+a)\ket{\Omega}.
\end{align*}

% PCT-SOURCE: pdf=127 print=115 kind=prose
The left-hand side has a continuation for complex $\Lambda$, providing a
continuation of the right-hand side by

% PCT-SOURCE: pdf=127 print=115 kind=display
\begin{align*}
 &\sum_{\alpha',\beta',\ldots,\gamma'}
 \bra{\Omega}\varphi_{\alpha'}(x_1)\psi_{\beta'}(x_2)\cdots
 \chi_{\gamma'}(x_n)\ket{\Omega}\\[-2pt]
 &\qquad\times S^{(\varphi)}_{\alpha\alpha'}(A,B)
 S^{(\psi)}_{\beta\beta'}(A,B)\cdots S^{(\chi)}_{\gamma\gamma'}(A,B)\\
 &\qquad=\bra{\Omega}\varphi_\alpha(\Lambda(A,B)x_1+a)
 \psi_\beta(\Lambda(A,B)x_2+a)\cdots
 \chi_\gamma(\Lambda(A,B)x_n+a)\ket{\Omega}.
\end{align*}

% PCT-SOURCE: pdf=127 print=115 kind=prose
This continuation defines the right-hand side for all points $\Lambda(A,B)x_j$
with $x_j-x_{j+1}\in\ExtendedTube_{n-1}$, as the $A,B$ vary over
$\mathrm{SL}(2,\C)\otimes\mathrm{SL}(2,\C)$. The continuation is
single-valued, by Theorem 2-11. Thus $W(\xi_1,\ldots,\xi_{n-1})$ is
continuable into $\ExtendedTube_{n-1}$.
\end{proof}

% PCT-SOURCE: pdf=127 print=115 kind=prose
An important special case is $\Lambda=-1$, which is connected in the complex
Lorentz group to $\Lambda=1$. By (1-27)
$S^{(\varphi)}(-1,1)=(-1)^j$, where $j$ is the number of undotted indices in
the field $\varphi$. Hence, if $J$ is the total number of undotted indices
in the fields appearing in the $W$, we get

% PCT-SOURCE: pdf=127 print=115 kind=display id=3-40
\begin{equation}
 W(\xi_1,\ldots,\xi_{n-1})
 =(-1)^J W(-\xi_1,-\xi_2,\ldots,-\xi_{n-1}),
 \tag{3-40}\label{eq:pct-3-40}
\end{equation}

% PCT-SOURCE: pdf=127 print=115 kind=prose
which holds at all points of holomorphy. It should be emphasized that
(3-40) holds, even if $P$, $C$, and $T$ do not define symmetries, since it
follows from invariance under $\mathcal{P}^{\uparrow}_{+}$ alone, combined
with the hypothesis about the mass spectrum of states.

% PCT-SOURCE: pdf=127 print=115 kind=prose
The final theorem relates the function $W$ to the holomorphic function
$W_\pi$, defined analogously for a permutation of the fields.

% PCT-SOURCE: pdf=127 print=115 kind=theorem id=3-6
\begin{theorem}
Suppose that

% PCT-SOURCE: pdf=127 print=115 kind=display
\begin{equation*}
 W(\xi_1,\xi_2,\ldots,\xi_{n-1})
 =\bra{\Omega}\varphi_1(x_1)\cdots\varphi_n(x_n)\ket{\Omega},
 \qquad \xi_j=x_j-x_{j+1},
\end{equation*}

% PCT-SOURCE: pdf=127 print=115 kind=display
\begin{equation*}
 W_\pi(\xi_1,\xi_2,\ldots,\xi_{n-1})
 =\bra{\Omega}\varphi_{i_1}(x_{i_1})\varphi_{i_2}(x_{i_2})\cdots
 \varphi_{i_n}(x_{i_n})\ket{\Omega},
\end{equation*}

% PCT-SOURCE: pdf=127 print=115 kind=prose
where $\pi$ is the permutation
$(1,2,\ldots,n)\longmapsto(i_1,i_2,\ldots,i_n)$; and suppose $W$ and
$W_\pi$ are the holomorphic functions given by the previous theorem. Then
$W$ and $W_\pi$ continue one another as one holomorphic function.
\end{theorem}

% PCT-SOURCE: pdf=127 print=115 kind=proof id=3-6
\begin{proof}
Suppose $x_1,x_2,\ldots,x_n$ are such that all the differences $x_i-x_j$ are
space-like.\footnote{Then we say that $(x_1,\ldots,x_n)$ is a \emph{totally
space-like point}.} Then $W$ and $W_\pi$ coincide in a real neighborhood of
this point, by local commutativity. We have only to show that these are

% PCT-SOURCE: pdf=128 print=116 kind=prose
points of holomorphy of both functions. This is achieved most simply by
applying the edge of the wedge theorem. Consider the distribution

% PCT-SOURCE: pdf=128 print=116 kind=display
\begin{equation*}
 \mathcal{W}'=\bra{\Omega}\varphi_n(x_n)\cdots\varphi_2(x_2)
 \varphi_1(x_1)\ket{\Omega}.
\end{equation*}

% PCT-SOURCE: pdf=128 print=116 kind=prose
Then the corresponding holomorphic function $W'$ is holomorphic if
$-\eta_j\in V_+$, by the previous theorem. Furthermore, the boundary values
$W'$ and $W$ agree at totally space-like points. Hence, by the edge of the
wedge theorem $W$ is holomorphic at such a point.
\end{proof}

% PCT-SOURCE: pdf=128 print=116 kind=prose
The theorem also follows directly from the fact that the extended tube and
the permuted extended tubes have a real environment in common (see Figure
2-4).

% PCT-SOURCE: pdf=128 print=116 kind=prose
The expectation values for a theory of a free field can be computed directly
from the definition of the field operators (3-14). For example, for a
hermitian scalar field the state $\ket{\Omega}$ is represented by
$(1,0,0,\ldots)$, the state $\varphi(f)\ket{\Omega}$ by
$\sqrt{\pi}(0,\widetilde f(p),0,\ldots)$. The scalar product is then

% PCT-SOURCE: pdf=128 print=116 kind=display
\begin{align*}
 \bra{\Omega}\varphi(f)\varphi(g)\ket{\Omega}
 &=\pi\int\widetilde f(p)\widetilde g(p)\,\dd\Omega_m(p)\\
 &=2\pi\int\widetilde f(p)\widetilde g(p)\theta(p^0)
 \delta(p^2+m^2)\,\dd^4p.
\end{align*}

% PCT-SOURCE: pdf=128 print=116 kind=prose
It follows that as a distribution

% PCT-SOURCE: pdf=128 print=116 kind=display
\begin{equation*}
 \bra{\Omega}\varphi(x)\varphi(y)\ket{\Omega}
 =\frac{1}{\ii}\Delta^+(x-y;m),
\end{equation*}

% PCT-SOURCE: pdf=128 print=116 kind=prose
where

% PCT-SOURCE: pdf=128 print=116 kind=display
\begin{equation*}
 \Delta^+(x;m)=\frac{\ii}{(2\pi)^3}
 \int\theta(p^0)\delta(p^2+m^2)\exp(+\ii p\!\cdot x)\,\dd^4p.
\end{equation*}

% PCT-SOURCE: pdf=128 print=116 kind=prose
A repeated application of (3-14) gives

% PCT-SOURCE: pdf=128 print=116 kind=display id=3-41
\begin{equation}
\begin{aligned}
 &\bra{\Omega}\varphi(x_1)\varphi(x_2)\cdots\varphi(x_n)\ket{\Omega}\\
 &\quad=
 \begin{cases}
 \begin{gathered}
  \displaystyle\sum_{\text{partitions}}
  \frac{1}{\ii}\Delta^+(x_{i_1}-x_{i_2})
  \frac{1}{\ii}\Delta^+(x_{i_3}-x_{i_4})\cdots\\[-2pt]
  \displaystyle\qquad\times\frac{1}{\ii}\Delta^+(x_{i_{n-1}}-x_{i_n})
 \end{gathered}
 &n\ \text{even},\\[6pt]
 0,&n\ \text{odd}.
 \end{cases}
\end{aligned}
\tag{3-41}\label{eq:pct-3-41}
\end{equation}

% PCT-SOURCE: pdf=128 print=116 kind=prose
where the sum is over all partitions of $n$ into $n/2$ disjoint two element
subsets $(i_1i_2)\allowbreak(i_3i_4)\allowbreak\cdots\allowbreak
(i_{n-1}i_n)$ with $i_{2k-1}<i_{2k}$. We leave
the proof to the reader. The study of the properties of the holomorphic
functions $W$ of a field theory has been called the \emph{linear program}.
An up-to-date summary of progress is to be found in Ref.~12. The idea is
that once we have computed the domain of holomorphy explicitly, we can
